% !TEX program = xelatex
\documentclass[a4paper,11pt]{article}

% ===================== Packages =====================
\usepackage{fontspec}
\usepackage{xeCJK}
\usepackage[a4paper,top=2.8cm,bottom=2.8cm,left=2.6cm,right=2.6cm,headheight=15pt,headsep=10pt,footskip=22pt]{geometry}
\usepackage{array}
\usepackage{amsmath}
\usepackage{amssymb}
\usepackage{booktabs}
\usepackage{tabularx}
\usepackage{longtable}
\usepackage{enumitem}
\usepackage{xcolor}
\usepackage{setspace}
\usepackage{titlesec}
\usepackage{titletoc}
\usepackage{float}
\usepackage{tikz}
\usepackage{fancyhdr}
\usepackage[font=small,labelfont=bf,labelsep=quad,skip=6pt]{caption}
\usepackage{listings}
\usepackage{hyperref}

\usetikzlibrary{arrows.meta,positioning,fit,calc,shapes.geometric,backgrounds}

% ===================== Fonts =====================
\IfFontExistsTF{Noto Sans CJK SC}{%
  \setCJKmainfont{Noto Sans CJK SC}%
  \setCJKsansfont{Noto Sans CJK SC}%
  \setCJKmonofont{Noto Sans CJK SC}%
}{%
  \IfFontExistsTF{PingFang SC}{%
    \setCJKmainfont{PingFang SC}%
    \setCJKsansfont{PingFang SC}%
    \setCJKmonofont{PingFang SC}%
  }{%
    \setCJKmainfont{AR PL UMing CN}%
    \setCJKsansfont{AR PL UKai CN}%
    \setCJKmonofont{AR PL UMing CN}%
  }%
}
\IfFontExistsTF{Noto Serif CJK SC}{%
  \setmainfont{Noto Serif CJK SC}%
}{%
  \IfFontExistsTF{Songti SC}{\setmainfont{Songti SC}}{\setmainfont{DejaVu Serif}}%
}
\IfFontExistsTF{Noto Sans CJK SC}{%
  \setsansfont{Noto Sans CJK SC}%
}{%
  \IfFontExistsTF{PingFang SC}{\setsansfont{PingFang SC}}{\setsansfont{DejaVu Sans}}%
}
\IfFontExistsTF{Noto Sans Mono CJK SC}{%
  \setmonofont{Noto Sans Mono CJK SC}%
}{%
  \IfFontExistsTF{Menlo}{\setmonofont{Menlo}}{\setmonofont{DejaVu Sans Mono}}%
}

% \XeTeXlinebreaklocale "zh"  % Tectonic's bundled ICU may not provide this locale
\XeTeXlinebreakskip = 0pt plus 1pt

% ===================== Hyperref =====================
\hypersetup{
  colorlinks=true,
  linkcolor=black,
  urlcolor=black,
  citecolor=black,
  pdftitle={ForL0 State Backend 编译构建指导书 v2.0},
  pdfauthor={ForL0 Project}
}

% ===================== Localization =====================
\renewcommand{\contentsname}{目\quad 录}
\renewcommand{\tablename}{表}
\renewcommand{\figurename}{图}

% ===================== Spacing =====================
\setstretch{1.28}
\setlength{\parindent}{2em}
\setlength{\parskip}{0.35em}
\setlength{\emergencystretch}{3em}
\sloppy
\setlist[itemize]{leftmargin=2em,itemsep=0.25em,topsep=0.35em,parsep=0pt}
\setlist[enumerate]{leftmargin=2.2em,itemsep=0.25em,topsep=0.35em,parsep=0pt}

% ===================== Headings =====================
\definecolor{accent}{RGB}{30,64,110}
\definecolor{rulegray}{RGB}{60,60,60}
\definecolor{softgray}{RGB}{245,246,248}
\definecolor{midgray}{RGB}{220,223,228}
\definecolor{linegray}{RGB}{120,124,132}
\definecolor{codebg}{RGB}{248,249,251}
\definecolor{codeborder}{RGB}{220,223,228}
\definecolor{codekw}{RGB}{30,64,110}
\definecolor{codecomment}{RGB}{90,100,110}
\definecolor{codestring}{RGB}{102,70,20}

\titleformat{\section}[hang]
  {\normalfont\large\bfseries\color{accent}}
  {\thesection}{0.8em}{}
  [\vspace{2pt}{\color{accent!60}\titlerule[0.6pt]}]
\titlespacing*{\section}{0pt}{1.6em}{0.9em}

\titleformat{\subsection}[hang]
  {\normalfont\normalsize\bfseries\color{accent!85!black}}
  {\thesubsection}{0.6em}{}
\titlespacing*{\subsection}{0pt}{1.2em}{0.55em}

\titleformat{\subsubsection}[hang]
  {\normalfont\normalsize\bfseries}
  {\thesubsubsection}{0.55em}{}
\titlespacing*{\subsubsection}{0pt}{0.9em}{0.4em}

% ===================== TOC styling =====================
\titlecontents{section}
  [0em]{\addvspace{0.55em}\bfseries}
  {\contentslabel{1.6em}}{}
  {\titlerule*[0.6pc]{.}\contentspage}
\titlecontents{subsection}
  [2.2em]{\addvspace{0.15em}}
  {\contentslabel{2.4em}}{}
  {\titlerule*[0.6pc]{.}\contentspage}
\titlecontents{subsubsection}
  [4.6em]{\addvspace{0.05em}\small}
  {\contentslabel{3.0em}}{}
  {\titlerule*[0.6pc]{.}\contentspage}

% ===================== Header / Footer =====================
\newcommand{\project}{ForL0 State Backend}
\pagestyle{fancy}
\fancyhf{}
\renewcommand{\headrulewidth}{0.4pt}
\renewcommand{\footrulewidth}{0pt}
\fancyhead[L]{\small\sffamily\color{accent}\project{} 编译构建指导书}
\fancyhead[R]{\small\sffamily\color{accent}v2.0}
\fancyfoot[C]{\small\thepage}

\fancypagestyle{plain}{%
  \fancyhf{}%
  \renewcommand{\headrulewidth}{0pt}%
  \fancyfoot[C]{\small\thepage}%
}

% ===================== Listings =====================
\lstdefinestyle{shellstyle}{
  basicstyle=\ttfamily\small,
  backgroundcolor=\color{codebg},
  frame=single,
  rulecolor=\color{codeborder},
  framesep=4pt,
  xleftmargin=6pt,
  xrightmargin=4pt,
  aboveskip=0.6em,
  belowskip=0.6em,
  breaklines=true,
  breakatwhitespace=false,
  showstringspaces=false,
  columns=fullflexible,
  keepspaces=true,
  tabsize=2,
  commentstyle=\color{codecomment}\itshape,
  stringstyle=\color{codestring},
  keywordstyle=\color{codekw}\bfseries,
  extendedchars=false,
  inputencoding=utf8,
  literate={\$}{{\$}}1 {\&}{{\&}}1 {\_}{{\_}}1
}
\lstdefinelanguage{shellcmd}{
  morekeywords={cd,ls,mkdir,cp,mv,rm,export,cmake,make,mvn,yum,sudo,wget,tar,ln,grep,nm,file,ldd,source,ln,tee,curl,echo,python3,pip3,bash,nproc,dlopen},
  morecomment=[l]{\#},
  morestring=[b]",
}
\lstdefinelanguage{yamlcmd}{
  morekeywords={state,backend,type,forl0,checkpoints,dir,env,java,opts,all,l0-cache,enabled,size,max-per-alloc,async-snapshots,execution,checkpointing,interval,mode},
  morecomment=[l]{\#},
  morestring=[b]",
}
\lstset{style=shellstyle,language=shellcmd}

% ===================== Table column types =====================
\newcolumntype{L}[1]{>{\raggedright\arraybackslash}p{#1}}
\newcolumntype{C}[1]{>{\centering\arraybackslash}p{#1}}

% ===================== Figure palette =====================
\definecolor{fgfill}{RGB}{248,249,251}
\definecolor{fgaccent}{RGB}{30,64,110}
\definecolor{fgaccentlight}{RGB}{223,231,242}
\definecolor{fgbox}{RGB}{255,255,255}
\definecolor{fgborder}{RGB}{90,95,105}
\definecolor{fgsubtle}{RGB}{235,237,240}
\definecolor{fgline}{RGB}{70,74,82}

\tikzset{
  figfont/.style={font=\footnotesize\sffamily},
  figsmall/.style={font=\scriptsize\sffamily,text=black!70},
  figtitle/.style={font=\footnotesize\sffamily\bfseries,text=fgaccent},
  box/.style={draw=fgborder, line width=0.5pt, rounded corners=2pt, fill=fgbox,
              minimum height=0.82cm, inner xsep=6pt, inner ysep=4pt, align=center, figfont},
  accentbox/.style={draw=fgaccent, line width=0.6pt, rounded corners=2pt, fill=fgaccentlight,
              minimum height=0.82cm, inner xsep=6pt, inner ysep=4pt, align=center, figfont},
  subtlebox/.style={draw=fgborder, line width=0.5pt, rounded corners=2pt, fill=fgsubtle,
              minimum height=0.82cm, inner xsep=6pt, inner ysep=4pt, align=center, figfont},
  lane/.style={draw=fgborder!55, line width=0.4pt, rounded corners=3pt, fill=fgfill},
  line/.style={-{Latex[length=2mm,width=1.6mm]}, draw=fgline, line width=0.55pt},
  thinline/.style={draw=fgline!80, line width=0.4pt}
}

% ===================== Title page macro =====================
\newcommand{\coverrule}{{\color{accent}\rule{\linewidth}{0.9pt}}}
\newcommand{\coverrulethin}{{\color{accent!55}\rule{\linewidth}{0.3pt}}}

\begin{document}

% =================================================
%                 COVER
% =================================================
\begin{titlepage}
\thispagestyle{empty}
\centering
\vspace*{5.5cm}

\coverrule
\vspace{2pt}
\coverrulethin

\vspace{1.6cm}
{\sffamily\huge\bfseries\color{accent} \project\par}
\vspace{0.8cm}
{\sffamily\LARGE\bfseries 编译构建指导书\par}

\vspace{1.6cm}
\coverrulethin
\vspace{2pt}
\coverrule

\vfill

{\sffamily\large\color{accent!85} 2026 年 8 月\par}

\vspace{1.2cm}
\end{titlepage}

% =================================================
%                 TOC
% =================================================
\pagenumbering{roman}
\pagestyle{plain}
\tableofcontents
\clearpage
\pagenumbering{arabic}
\setcounter{page}{1}
\pagestyle{fancy}

% =================================================
%                 CONTENT
% =================================================
\section{文档概述}

本文档给出 \project{} 从基础环境准备、原生引擎编译、Java 主工程打包到基线 Benchmark 构建的完整流程，仅覆盖编译构建相关内容。集群部署、作业配置、运行期验证等使用层面内容由《使用说明书》单独说明；架构与数据结构设计由《设计说明书》给出。

\subsection{目标平台}

文中提到的所有路径、依赖包和命令均假定在如下环境下执行：

\begin{table}[H]
\centering
\small
\begin{tabularx}{\textwidth}{L{3.2cm} X}
\toprule
\textbf{项目} & \textbf{要求} \\
\midrule
处理器 & 鲲鹏系列处理器 \\
操作系统 & openEuler \\
架构 & aarch64（ARM64） \\
内核特性 & 支持鲲鹏 L0 Cache 的内核驱动及 \texttt{/dev/hisi\_l0} 或 \texttt{/dev/l0} 设备节点（若计划启用 L0 HotCache） \\
\bottomrule
\end{tabularx}
\end{table}

L0 Cache 相关的硬件与驱动并非编译期依赖。若目标机器缺少 L0 硬件或内存池库，L0 HotCache 保持关闭，权威状态仍由原生 SwissTable 正常提供。

\subsection{构建产物}

按照本文档的构建流程，最终会得到下列产物：

\begin{table}[H]
\centering
\small
\begin{tabularx}{\textwidth}{L{5.8cm} X}
\toprule
\textbf{产物} & \textbf{说明} \\
\midrule
\texttt{libforl0\_engine.so} & 原生状态引擎共享库（aarch64 版本） \\
\nolinkurl{flink-statebackend-forL0-1.0-SNAPSHOT.jar} & Java 控制层与 JNI 桥接打包的状态后端 JAR \\
\nolinkurl{wordcount-benchmark-*.jar} & WordCount 基线性能测试 JAR \\
\nolinkurl{nexmark-flink-0.3-SNAPSHOT.jar} & NexMark 基线性能测试 JAR \\
\bottomrule
\end{tabularx}
\end{table}

\section{基础环境准备}

基础环境准备包括系统构建工具、JDK、Maven、CMake、Flink 运行环境，以及 L0 Cache 相关的可选依赖。若构建机与部署机分离，基础构建工具可只安装在构建机上；JDK、Flink 与 L0 库一般需要部署到运行机。

\subsection{系统构建工具}

C++ 原生引擎采用 CMake + Make 构建，对 GCC 版本有一定要求。在 openEuler 上可通过包管理器直接安装：

\begin{lstlisting}[language=shellcmd]
sudo yum install -y gcc gcc-c++ cmake3 make

# 若默认仅有 cmake3 可执行文件，建议创建 cmake 软链接
sudo ln -sf /usr/bin/cmake3 /usr/bin/cmake

gcc --version   # 要求 >= 7.0，支持 C++17
\end{lstlisting}

部分发行版自带的 CMake 版本低于 3.14，无法满足原生引擎构建。此时可使用官方 aarch64 预编译包：

\begin{lstlisting}[language=shellcmd]
wget https://github.com/Kitware/CMake/releases/download/v3.28.0/cmake-3.28.0-linux-aarch64.tar.gz
tar xzf cmake-3.28.0-linux-aarch64.tar.gz -C /opt/
export PATH=/opt/cmake-3.28.0-linux-aarch64/bin:$PATH
cmake --version
\end{lstlisting}

\subsection{JDK 与 Maven}

Java 主工程使用 JDK 1.8 与 Maven 3.6 及以上版本。在鲲鹏平台建议使用华为毕昇 JDK，也可直接使用 OpenJDK 提供的 aarch64 包：

\begin{lstlisting}[language=shellcmd]
# OpenJDK（系统源）
sudo yum install -y java-1.8.0-openjdk-devel

# 或使用 BiSheng JDK（鲲鹏优化，从华为镜像站下载 aarch64 版本后解压至 /opt）
export JAVA_HOME=/opt/bisheng-jdk1.8.0
export PATH=$JAVA_HOME/bin:$PATH

java -version
javac -version
\end{lstlisting}

Maven 使用官方二进制包安装：

\begin{lstlisting}[language=shellcmd]
wget https://archive.apache.org/dist/maven/maven-3/3.8.8/binaries/apache-maven-3.8.8-bin.tar.gz
tar xzf apache-maven-3.8.8-bin.tar.gz -C /opt/
export MAVEN_HOME=/opt/apache-maven-3.8.8
export PATH=$MAVEN_HOME/bin:$PATH
mvn --version
\end{lstlisting}

\subsection{Apache Flink 运行环境}

部署机需要具备 Apache Flink 1.20.3 二进制发行版。Flink 自身为 Scala 2.12 构建，无架构差异，可直接从官方下载：

\begin{lstlisting}[language=shellcmd]
wget https://archive.apache.org/dist/flink/flink-1.20.3/flink-1.20.3-bin-scala_2.12.tgz
tar xzf flink-1.20.3-bin-scala_2.12.tgz -C /opt/
export FLINK_HOME=/opt/flink-1.20.3
$FLINK_HOME/bin/flink --version
\end{lstlisting}

\subsection{L0 Cache 依赖检测}

鲲鹏 L0 Cache 加速依赖 L0 设备节点和内存池共享库 \texttt{libl0mempool.so}。二者均为运行期依赖，构建期无需参与链接，但在部署前应提前检查。

\begin{lstlisting}[language=shellcmd]
# 检查设备节点、共享库与可调用符号
bash docker/lib/l0_detector.sh

# 也可分别检查设备节点
ls -la /dev/hisi_l0 /dev/l0 2>/dev/null

# 检查共享库解析结果
ldconfig -p 2>/dev/null | grep libl0mempool || true
\end{lstlisting}

设备节点、共享库、必要符号及首次 L0 分配均通过后，运行时才会激活 HotCache。任一条件不满足时，非 strict 配置保持 HotCache 关闭，SwissTable 路径不受影响。构建期无需任何 L0 依赖参与。

\subsection{环境变量汇总}

为便于后续命令复用，建议将关键变量固化到 \texttt{\~{}/.bashrc}：

\begin{lstlisting}[language=shellcmd]
export JAVA_HOME=/opt/bisheng-jdk1.8.0
export MAVEN_HOME=/opt/apache-maven-3.8.8
export FLINK_HOME=/opt/flink-1.20.3
export PATH=$JAVA_HOME/bin:$MAVEN_HOME/bin:$FLINK_HOME/bin:$PATH
\end{lstlisting}

\section{原生引擎构建}

原生引擎是 \project{} 的核心组件，包含 SwissTable 哈希存储、状态引擎、Copy-On-Write 快照实现以及 JNI 桥接层。原生引擎须先于 Java 主工程构建，Java 侧的测试与运行均依赖其产物。

\subsection{构建步骤}

原生引擎位于 \texttt{src/main/native} 下。交付构建与离线脚本使用仓库内 Makefile，推荐流程如下：

\begin{lstlisting}[language=shellcmd]
cd src/main/native
make clean
make -j$(nproc)
make install
\end{lstlisting}

\texttt{make install} 会将生成的 \texttt{libforl0\_engine.so} 复制到 \texttt{src/main/resources/native/}，随后 Java 测试和打包流程即可引用该产物。CMake 主要用于构建原生单元测试，也会在共享库构建完成后复制同一产物。

\subsection{完整重建与产物一致性}

native 引擎的大量实现位于模板头文件中。头文件、JNI 声明或类型布局发生变化后，应执行 \texttt{make clean} 再重新构建，避免残留 object 与当前源码不一致。

部署前必须保持以下副本一致：\nolinkurl{src/main/resources/native/libforl0_engine.so}、\nolinkurl{docker/deploy/libforl0_engine.so}、backend JAR 内嵌 native resource，以及 Flink \texttt{lib/}/\texttt{native/} 中的运行副本。仅替换 JAR 或仅替换外部 \texttt{.so} 都可能产生 Java/JNI/native 版本错配。

\subsection{鲲鹏平台默认启用的编译选项}

在 aarch64 上，Makefile 与 CMake 脚本会根据平台自动启用下列优化，不需要在命令行显式指定：

\begin{itemize}
  \item \texttt{-O3 -DNDEBUG -march=native}：以鲲鹏 CPU 为目标的最高级别优化。
  \item \texttt{-DFORL0\_NEON=1}：启用 NEON 路径的 SWAR 并行匹配实现，用于 SwissTable 的控制字节比较。
  \item \texttt{\$\{CMAKE\_DL\_LIBS\}} 链接：仅链接 \texttt{libdl}，以便运行时通过 \texttt{dlopen} 动态加载 L0 内存池库。
\end{itemize}

构建产物应为 aarch64 架构的 ELF 共享库，可通过如下命令核对：

\begin{lstlisting}[language=shellcmd]
file src/main/resources/native/libforl0_engine.so
# 预期: ELF 64-bit LSB shared object, ARM aarch64, ...
nm -D src/main/resources/native/libforl0_engine.so | grep Java_org_apache | head
\end{lstlisting}

\subsection{可选的 C++ 单元测试}

原生单元测试由 CMake 构建，\texttt{FORL0\_BUILD\_TESTS} 默认启用：

\begin{lstlisting}[language=shellcmd]
cd src/main/native
rm -rf build && mkdir build && cd build
cmake .. -DCMAKE_BUILD_TYPE=Release -DFORL0_BUILD_TESTS=ON
make -j$(nproc)
ctest --output-on-failure
\end{lstlisting}

\subsection{CMake 选项汇总}

\begin{table}[H]
\centering
\small
\begin{tabularx}{\textwidth}{L{4.2cm} L{2.4cm} X}
\toprule
\textbf{选项} & \textbf{默认值} & \textbf{说明} \\
\midrule
\texttt{CMAKE\_BUILD\_TYPE} & 未指定 & 未指定时使用调试编译参数；生产交付必须显式设置为 \texttt{Release}。 \\
\texttt{FORL0\_BUILD\_TESTS} & \texttt{ON} & 置为 \texttt{OFF} 时跳过 C++ 单元测试可执行文件。 \\
\bottomrule
\end{tabularx}
\end{table}

% ---------- 图 1: 构建流水线 ----------
\begin{figure}[htbp]
\centering
\begin{tikzpicture}[x=1cm,y=1cm]
  \node[lane, minimum width=14.6cm, minimum height=3.6cm] at (7.3,1.8) {};

  \node[box, minimum width=2.4cm] (env)    at (1.9,1.8) {环境准备\\JDK / Maven\\CMake / GCC};
  \node[accentbox, minimum width=2.4cm] (cmake)  at (5.2,1.8) {Native 构建\\Release + NEON};
  \node[accentbox, minimum width=2.4cm] (make)   at (8.5,1.8) {make -j\\\texttt{libforl0\_engine.so}};
  \node[box, minimum width=2.4cm] (mvn)    at (11.8,1.8) {Maven 打包\\ForL0 JAR};

  \draw[line] (env.east)    -- (cmake.west);
  \draw[line] (cmake.east)  -- (make.west);
  \draw[line] (make.east)   -- (mvn.west);
\end{tikzpicture}
\caption{\project{} 构建流水线}
\label{fig:build-pipeline}
\end{figure}

\section{Java 主工程构建}

原生引擎产物就绪后即可构建 Java 主工程。主工程负责 Flink 侧后端适配、JNI 调用入口以及测试用例打包。

\subsection{编译与单元测试}

基础编译与测试命令如下：

\begin{lstlisting}[language=shellcmd]
# 回到项目根目录
cd $PROJECT_ROOT

# 仅编译源码
mvn clean compile

# 编译 + 运行单元测试
mvn test
\end{lstlisting}

Maven Surefire 已预置 \texttt{-Djava.library.path=src/main/resources/native}，测试阶段会自动加载原生库，无需额外配置。

\subsection{打包 JAR}

构建最终交付 JAR：

\begin{lstlisting}[language=shellcmd]
mvn package -DskipTests
ls target/flink-statebackend-forL0-1.0-SNAPSHOT.jar
\end{lstlisting}

生产交付建议使用带测试的完整构建：

\begin{lstlisting}[language=shellcmd]
mvn clean package
\end{lstlisting}

若构建机无法访问公共 Maven 仓库，需在 \texttt{\textasciitilde/.m2/settings.xml} 中配置内网镜像，或先行准备离线仓库。

\section{Benchmark 构建}

\project{} 交付物中包含两套基线性能测试：WordCount 基线和 NexMark 基线。它们与主工程解耦，可独立构建，也可在部署后单独提交至 Flink 集群。

\subsection{WordCount Benchmark}

\begin{lstlisting}[language=shellcmd]
cd benchmark/wordcount
mvn clean package -DskipTests
ls target/wordcount-benchmark-*.jar
\end{lstlisting}

\subsection{NexMark Benchmark}

\begin{lstlisting}[language=shellcmd]
cd benchmark/nexmark-src
mvn clean package -DskipTests
ls nexmark-flink/target/nexmark-flink-0.3-SNAPSHOT.jar
ls nexmark-flink/target/nexmark-flink-bin/nexmark-flink/bin
\end{lstlisting}

\subsection{（可选）JMH Benchmark}

Flink 官方的 JMH Benchmark 用于回归对比，可按需构建：

\begin{lstlisting}[language=shellcmd]
cd benchmark/flink-benchmarks
./mvnw clean package -DskipTests
\end{lstlisting}

\subsection{Benchmark 脚本依赖}

Benchmark 运行脚本使用 Python 编写。运行环境需安装 Python 3.8 及以上版本：

\begin{lstlisting}[language=shellcmd]
sudo yum install -y python3 python3-pip
cd benchmark/
pip3 install -r requirements.txt

# 封闭网络环境下使用随源码交付的离线安装包
pip3 install --no-index --find-links=offline-packages/ -r requirements.txt
\end{lstlisting}

\section{一键构建与离线打包}

仓库提供正式的一键构建与离线打包入口。该脚本依次构建后端 JAR、native 库、Benchmark 产物，并收集离线运行依赖：

\begin{lstlisting}[language=shellcmd]
cd $PROJECT_ROOT
chmod +x docker/package_offline_bundle.sh

# 输出到指定目录，并固定目标 Python ABI
./docker/package_offline_bundle.sh \
  --arch arm64 \
  --python-version 3.10 \
  --output-dir "$PWD/offline-artifacts/forl0-offline-linux-arm64-py310"
\end{lstlisting}

\section{常见问题与排查}

\begin{table}[H]
\centering
\small
\begin{tabularx}{\textwidth}{L{4.8cm} L{3.2cm} X}
\toprule
\textbf{现象} & \textbf{可能原因} & \textbf{处理方式} \\
\midrule
\texttt{CMake Error: Could not find JNI} & \texttt{JAVA\_HOME} 未正确设置 & \nolinkurl{export JAVA_HOME=/opt/bisheng-jdk1.8.0} 后重新配置 \\
\texttt{unrecognized option `-march=native'} & GCC 版本低于 7.0 & 升级 GCC（\texttt{sudo yum install devtoolset-8}）或使用 BiSheng 编译器 \\
\texttt{jni.h: No such file or directory} & 缺少 JDK 开发包 & 安装 \nolinkurl{java-1.8.0-openjdk-devel} 或确认 BiSheng JDK 路径完整 \\
\texttt{cannot find -lstdc++fs} & libstdc++ 版本过低 & 安装 \texttt{libstdc++-devel} 或切换至新版 GCC \\
\texttt{mvn test} 抛出 \texttt{UnsatisfiedLinkError} & 原生库缺失或与 Java 代码版本不一致 & 重新执行 \S3 节的 \texttt{make clean}、\texttt{make} 和 \texttt{make install} \\
 NexMark 编译卡在依赖下载 & 仓库不可达 & 配置内网 Maven 镜像或使用离线仓库 \\
\bottomrule
\end{tabularx}
\end{table}

\section{附录}

\subsection{项目目录结构}

以下目录结构摘录了本文档涉及的构建相关路径，其他内容参见项目源码：

\begin{lstlisting}[language=shellcmd]
forL0-state-backend/
|-- pom.xml                          # Java 主工程 POM
|-- src/
|   |-- main/
|   |   |-- java/                    # Java 源码（Flink 适配与 JNI 入口）
|   |   |   `-- org/apache/flink/state/forl0/
|   |   |-- native/                  # 原生引擎源码
|   |   |   |-- CMakeLists.txt
|   |   |   |-- Makefile
|   |   |   |-- engine/              # 引擎与 HotCache 实现
|   |   |   |   |-- swiss_table.h    # Swiss Tables 哈希表
|   |   |   |   |-- allocator.h      # 分配器抽象
|   |   |   |   |-- hot_cache.h       # L0 HotCache
|   |   |   |   `-- state_engine.h
|   |   |   |-- jni/                 # JNI 桥接
|   |   |   |-- checkpoint/          # 快照序列化
|   |   |   `-- test/                # C++ 单元测试
|   |   `-- resources/native/         # native 构建复制的共享库
|   `-- test/java/                   # Java 测试代码
|-- benchmark/
|   |-- wordcount/                   # WordCount Benchmark
|   |-- nexmark-src/                 # NexMark Benchmark
|   |-- flink-benchmarks/            # JMH Benchmark
|   `-- scripts/                     # Python 运行脚本
`-- target/                          # Maven 产物
\end{lstlisting}

\end{document}
